% GENERATED FILE. Edit canonical lessons, then run npm run generate:books.
% canonical-source-hash: fnv1a64:77e1a8f197854bf7
% canonical-lessons: JA-C185-waribiki, JA-C185-yasuuri, JA-C185-urikire, JA-C185-kakaru, JA-C185-tariru

\chapter{Sale, Discount and Sold Out}
\label{ch:ja-sale-discount-and-sold-out}
\hlchaptermodality{185}

\begin{chapteropening}
\textbf{By the end of this chapter:} \emph{I can read sale, discount and sold-out signs.}

Complete the last lesson of chapter 185: I can read sale, discount and sold-out signs.
\end{chapteropening}

\section[\texorpdfstring{\ja{わりびき} (waribiki)}{わりびき (waribiki)}]{\ja{わりびき} (waribiki) — discount; reduction}

\label{lesson:JA-C185-waribiki}



\begin{quote}
Before the new one: say the Japanese for all of it, then the Japanese for exactly.
\end{quote}

\subsection*{You'll want to know: \ja{わりびき}}
\textbf{\ja{わりびき}} — \emph{waribiki} — \textquotedblleft{}discount; reduction\textquotedblright{}.

\begin{quote}
\textbf{\ja{この} \ja{くつは} \ja{わりびきです。}} — \emph{kono kutsu wa waribiki desu} — \textquotedblleft{}These shoes are discounted.\textquotedblright{}
\end{quote}

\begin{cousinweb}
\textbf{\ja{わりびき}} (\emph{waribiki}) is native Japanese, built from \emph{wari}, the stem of \emph{waru}, \textquotedblleft{}to divide\textquotedblright{}, and \emph{biki}, from \emph{hiku}, \textquotedblleft{}to pull, to take off\textquotedblright{} (JMdict, KANJIDIC2).
\end{cousinweb}

\begin{grammarlens}[title={What you've built}]
One more word for everyday life at A2.
\end{grammarlens}

\subsection*{Your turn}
\begin{itemize}
\raggedright
  \item \emph{Say it:} \emph{waribiki}
  \item \emph{Say it:} \emph{waribiki}, once more
  \item \emph{Say it:} \emph{kono kutsu wa waribiki desu} — These shoes are discounted.
  \item \emph{Recall:} say the Japanese for all of it, then the Japanese for exactly, then say \emph{waribiki} again
\end{itemize}

\begin{tcolorbox}[breakable,colback=teal!4,colframe=teal!35!black,title={Before you move on}]
What does \ja{わりびき} mean? (\textquotedblleft{}Discount; reduction\textquotedblright{}.) Say it once more.
\end{tcolorbox}
\section[\texorpdfstring{\ja{やすうり} (yasuuri)}{やすうり (yasuuri)}]{\ja{やすうり} (yasuuri) — bargain sale}

\label{lesson:JA-C185-yasuuri}



\begin{quote}
Before the new one: say the Japanese for exactly, then the Japanese for discount.
\end{quote}

\subsection*{You'll want to know: \ja{やすうり}}
\textbf{\ja{やすうり}} — \emph{yasuuri} — \textquotedblleft{}bargain sale\textquotedblright{}.

\begin{quote}
\textbf{\ja{しゅうまつは} \ja{やすうりが} \ja{あります。}} — \emph{shūmatsu wa yasuuri ga arimasu} — \textquotedblleft{}There is a bargain sale at the weekend.\textquotedblright{}
\end{quote}

\begin{cousinweb}
\textbf{\ja{やすうり}} (\emph{yasuuri}) is native Japanese, built from \emph{yasu}, the stem of \emph{yasui}, \textquotedblleft{}cheap\textquotedblright{}, and \emph{uri}, \textquotedblleft{}selling\textquotedblright{}, from \emph{uru}, \textquotedblleft{}to sell\textquotedblright{} (JMdict, KANJIDIC2).
\end{cousinweb}

\begin{grammarlens}[title={What you've built}]
One more word for everyday life at A2.
\end{grammarlens}

\subsection*{Your turn}
\begin{itemize}
\raggedright
  \item \emph{Say it:} \emph{yasuuri}
  \item \emph{Say it:} \emph{yasuuri}, once more
  \item \emph{Say it:} \emph{shūmatsu wa yasuuri ga arimasu} — There is a bargain sale at the weekend.
  \item \emph{Recall:} say the Japanese for exactly, then the Japanese for discount, then say \emph{yasuuri} again
\end{itemize}

\begin{tcolorbox}[breakable,colback=teal!4,colframe=teal!35!black,title={Before you move on}]
What does \ja{やすうり} mean? (\textquotedblleft{}Bargain sale\textquotedblright{}.) Say it once more.
\end{tcolorbox}
\section[\texorpdfstring{\ja{うりきれ} (urikire)}{うりきれ (urikire)}]{\ja{うりきれ} (urikire) — sold out}

\label{lesson:JA-C185-urikire}



\begin{quote}
Before the new one: say the Japanese for discount, then the Japanese for bargain sale.
\end{quote}

\subsection*{You'll want to know: \ja{うりきれ}}
\textbf{\ja{うりきれ}} — \emph{urikire} — \textquotedblleft{}sold out\textquotedblright{}.

\begin{quote}
\textbf{\ja{やすい} \ja{たまごは} \ja{もう} \ja{うりきれです。}} — \emph{yasui tamago wa mō urikire desu} — \textquotedblleft{}The cheap eggs are already sold out.\textquotedblright{}
\end{quote}

\begin{cousinweb}
\textbf{\ja{うりきれ}} (\emph{urikire}) is native Japanese, built from \emph{uri}, \textquotedblleft{}selling\textquotedblright{}, and \emph{kire}, \textquotedblleft{}running out\textquotedblright{}, from \emph{kireru}, a verb on the same stem as \emph{kiru}, \textquotedblleft{}to cut\textquotedblright{} (JMdict, KANJIDIC2).
\end{cousinweb}

\begin{grammarlens}[title={What you've built}]
One more word for everyday life at A2.
\end{grammarlens}

\subsection*{Your turn}
\begin{itemize}
\raggedright
  \item \emph{Say it:} \emph{urikire}
  \item \emph{Say it:} \emph{urikire}, once more
  \item \emph{Say it:} \emph{yasui tamago wa mō urikire desu} — The cheap eggs are already sold out.
  \item \emph{Recall:} say the Japanese for discount, then the Japanese for bargain sale, then say \emph{urikire} again
\end{itemize}

\begin{tcolorbox}[breakable,colback=teal!4,colframe=teal!35!black,title={Before you move on}]
What does \ja{うりきれ} mean? (\textquotedblleft{}Sold out\textquotedblright{}.) Say it once more.
\end{tcolorbox}
\section[\texorpdfstring{\ja{かかる} (kakaru)}{かかる (kakaru)}]{\ja{かかる} (kakaru) — to cost; to take (time), to hang}

\label{lesson:JA-C185-kakaru}



\begin{quote}
Before the new one: say the Japanese for bargain sale, then the Japanese for sold out.
\end{quote}

\subsection*{You'll want to know: \ja{かかる}}
\textbf{\ja{かかる}} — \emph{kakaru} — \textquotedblleft{}to cost; to take (time), to hang\textquotedblright{}.

\begin{quote}
\textbf{\ja{くるまは} \ja{おかねが} \ja{かかります。}} — \emph{kuruma wa okane ga kakarimasu} — \textquotedblleft{}A car costs money.\textquotedblright{}
\end{quote}

Polite \textbf{\ja{かかります}}, \emph{kakarimasu}; a godan verb.

\begin{cousinweb}
\textbf{\ja{かかる}} (\emph{kakaru}) is a native Japanese word (wago), inherited rather than borrowed from Chinese; its kanji was attached later as a kun'yomi (JMdict, KANJIDIC2).
\end{cousinweb}

\begin{grammarlens}[title={What you've built}]
One more word for everyday life at A2.
\end{grammarlens}

\subsection*{Your turn}
\begin{itemize}
\raggedright
  \item \emph{Say it:} \emph{kakaru}
  \item \emph{Say it:} \emph{kakaru}, once more
  \item \emph{Say it:} \emph{kuruma wa okane ga kakarimasu} — A car costs money.
  \item \emph{Recall:} say the Japanese for bargain sale, then the Japanese for sold out, then say \emph{kakaru} again
\end{itemize}

\begin{tcolorbox}[breakable,colback=teal!4,colframe=teal!35!black,title={Before you move on}]
What does \ja{かかる} mean? (\textquotedblleft{}To cost; to take (time), to hang\textquotedblright{}.) Say it once more.
\end{tcolorbox}
\section[\texorpdfstring{\ja{たりる} (tariru)}{たりる (tariru)}]{\ja{たりる} (tariru) — to be enough; to suffice}

\label{lesson:JA-C185-tariru}



\begin{quote}
Before the new one: say the Japanese for sold out, then the Japanese for to cost.
\end{quote}

\subsection*{You'll want to know: \ja{たりる}}
\textbf{\ja{たりる}} — \emph{tariru} — \textquotedblleft{}to be enough; to suffice\textquotedblright{}.

\begin{quote}
\textbf{\ja{おかねが} \ja{たりません。}} — \emph{okane ga tarimasen} — \textquotedblleft{}There is not enough money.\textquotedblright{}
\end{quote}

Polite \textbf{\ja{たります}}, \emph{tarimasu}; an ichidan verb.

\begin{cousinweb}
\textbf{\ja{たりる}} (\emph{tariru}) is a native Japanese word (wago), inherited rather than borrowed from Chinese; its kanji was attached later as a kun'yomi (JMdict, KANJIDIC2).
\end{cousinweb}

\begin{grammarlens}[title={What you've built}]
One more word for everyday life at A2.
\end{grammarlens}

\subsection*{Your turn}
\begin{itemize}
\raggedright
  \item \emph{Say it:} \emph{tariru}
  \item \emph{Say it:} \emph{tariru}, once more
  \item \emph{Say it:} \emph{okane ga tarimasen} — There is not enough money.
  \item \emph{Recall:} say the Japanese for sold out, then the Japanese for to cost, then say \emph{tariru} again
\end{itemize}

\begin{tcolorbox}[breakable,colback=teal!4,colframe=teal!35!black,title={Before you move on}]
What does \ja{たりる} mean? (\textquotedblleft{}To be enough; to suffice\textquotedblright{}.) Say it once more.
\end{tcolorbox}
